\documentclass[12pt, twoside]{article}
%\documentclass[12pt, twoside]{article}
\usepackage[letterpaper, margin=1in, headsep=0.2in]{geometry}
\setlength{\headheight}{0.6in}
%\usepackage[english]{babel}
\usepackage[utf8]{inputenc}
\usepackage{microtype}
\usepackage{amsmath}
\usepackage{amssymb}
%\usepackage{amsfonts}
\usepackage[nomessages]{fp} %\FPeval{\var-name}{2*sin(pi/6)}
\usepackage{siunitx} %units in math. eg 20\milli\meter
\usepackage{yhmath} % for arcs, overparenth command
\usepackage{tikz} %graphics
\usetikzlibrary{quotes, angles, arrows, arrows.meta}
\usepackage{pgfplots}
\usepgfplotslibrary{fillbetween}
\pgfplotsset{compat=1.16}
\usepackage{graphicx} %consider setting \graphicspath{{images/}}
\usepackage{parskip} %no paragraph indent
\usepackage{enumitem}
\usepackage{multicol}
\usepackage{venndiagram}

\usepackage{fancyhdr}
\pagestyle{fancy}
\fancyhf{}
\renewcommand{\headrulewidth}{0pt} % disable the underline of the header
\raggedbottom
\hfuzz=2mm %suppresses overfull box warnings

\usepackage{hyperref}
\setlength{\parskip}{0.3\baselineskip}

\fancyhead[LE]{\thepage}
\fancyhead[RO]{Markscheme}
\fancyhead[LO]{La Scuola d'Italia / Dr.\ Huson / IB Math 12 \\* 29 September 2026}

\begin{document}
\subsubsection*{1.5 Classwork: Box-and-Whisker Plots --- Markscheme}

\emph{Calculator allowed. Marks are not printed on the student sheet; they are
recorded here. Total: 19 marks.}

\emph{\textbf{Throughout:} a whisker runs to the most extreme value that is
\emph{not} an outlier, and each outlier is drawn as a separate cross. The most
likely error is a whisker drawn all the way to the outlier (or to the fence
value $Q_{3} + 1.5 \times \text{IQR}$, which is not a data value).}

\begin{enumerate}[itemsep=0.3cm]

%--- Problem 1: Five-number summary from raw data; IQR; outlier test; draw the box plot [9 marks] ---
\item[1.] [Maximum mark: 9]

Ordered: $12, 15, 18, 20, 21, 22, 24, \mathbf{25}, 26, 28, 30, 31, 33, 35,
58$ \quad ($n = 15$).

\emph{\textbf{GDC:} enter the fifteen times as a list and run
\texttt{One-Variable Statistics}. The TI-Nspire excludes the median from each
half, so $Q_{1}$ is the median of the lowest seven values and $Q_{3}$ the
median of the highest seven.}

\textbf{(a)} Median $= 25$ \hfill \textbf{A1}

$Q_{1} = 20$, \quad $Q_{3} = 31$ \hfill \textbf{A1}

\emph{Both quartiles are needed for the second A1. Accept $Q_{1} = 20.5$ and
$Q_{3} = 30.5$ from a hand method that includes the median in each half;
feedback: the GDC convention is the one to use.}

\textbf{(b)} $\text{IQR} = 31 - 20 = 11$ \hfill \textbf{A1}

\emph{FT: award A1 for $Q_{3} - Q_{1}$ from the candidate's own quartiles.
The range, $58 - 12 = 46$, earns no mark.}

\textbf{(c)} $Q_{3} + 1.5 \times \text{IQR} = 31 + 1.5 \times 11$
\hfill \textbf{M1}

$= 47.5$ \hfill \textbf{A1}

$58 > 47.5$, so $58$ is an outlier. \hfill \textbf{R1}

\emph{M1 for $1.5 \times$ their IQR added to their $Q_{3}$. R1 for the
comparison with $58$ stated, not just the fence value. Adding
$1.5 \times \text{IQR}$ to the median or the mean earns no marks. The lower
fence, $20 - 16.5 = 3.5$, is not required. FT: award all three marks from the
candidate's own $Q_{3}$ and IQR, provided their fence is below $58$.}

\textbf{(d)}

\begin{center}
\begin{tikzpicture}
  \begin{scope}[x=0.2cm]
    \draw (0,0)--(60,0);
    \foreach \t in {0,5,...,60} \draw (\t,-0.12)--(\t,0.12);
    \foreach \t in {0,10,...,60} \node[below] at (\t,-0.12) {\footnotesize \t};
    \draw[thick] (20,0.6) rectangle (31,1.4);
    \draw[thick] (25,0.6)--(25,1.4);
    \draw[thick] (12,1.0)--(20,1.0);
    \draw[thick] (31,1.0)--(35,1.0);
    \draw[thick] (12,0.8)--(12,1.2);
    \draw[thick] (35,0.8)--(35,1.2);
    \node at (58,1.0) {$\times$};
  \end{scope}
\end{tikzpicture}
\end{center}

Box from $Q_{1} = 20$ to $Q_{3} = 31$ with the median line at $25$
\hfill \textbf{A1}

Whiskers to $12$ and $35$ \hfill \textbf{A1}

Outlier $58$ marked with a cross \hfill \textbf{A1}

\emph{Award each mark independently; plotted values within half a small
division ($\pm 1$ minute) are accepted. A whisker drawn to $58$ or to $47.5$
loses the second A1. FT: award the first A1 for a box drawn at the
candidate's own quartiles from (a).}

\newpage

%--- Problem 2: Read a box plot; outlier test from a five-number summary; draw and compare [10 marks] ---
\item[2.] [Maximum mark: 10]

Class A (from the plot): minimum $40$, $Q_{1} = 52$, median $60$,
$Q_{3} = 68$, maximum $84$.

\textbf{(a)} Median $= 60$ \hfill \textbf{A1}

$\text{IQR} = 68 - 52 = 16$ \hfill \textbf{A1}

\emph{The marks are independent. The range, $84 - 40 = 44$, earns no mark for
the IQR.}

\textbf{(b)} Class B: $\text{IQR} = 76 - 56 = 20$, \quad
$1.5 \times 20 = 30$ \hfill \textbf{M1}

Lower fence $56 - 30 = 26$, \quad upper fence $76 + 30 = 106$
\hfill \textbf{A1}

$22 < 26$, so $22$ is an outlier; $94 < 106$, so $94$ is not.
\hfill \textbf{R1}

\emph{M1 for $1.5 \times$ the Class B IQR. A1 for both fences. R1 for both
conclusions, each linked to its fence. Using the Class A IQR earns no
marks.}

\textbf{(c)}

\begin{center}
\begin{tikzpicture}
  \node[left] at (0,1.9) {\footnotesize Class A};
  \node[left] at (0,0.8) {\footnotesize Class B};
  \begin{scope}[xshift=-2.8cm, x=0.14cm]
    \draw (20,0)--(100,0);
    \foreach \s in {20,25,...,100} \draw (\s,-0.12)--(\s,0.12);
    \foreach \s in {20,30,...,100} \node[below] at (\s,-0.12) {\footnotesize \s};
    % Class A: 40, 52, 60, 68, 84
    \draw[thick] (52,1.6) rectangle (68,2.2);
    \draw[thick] (60,1.6)--(60,2.2);
    \draw[thick] (40,1.9)--(52,1.9);
    \draw[thick] (68,1.9)--(84,1.9);
    \draw[thick] (40,1.75)--(40,2.05);
    \draw[thick] (84,1.75)--(84,2.05);
    % Class B: outlier 22; 41, 56, 66, 76, 94
    \draw[thick] (56,0.5) rectangle (76,1.1);
    \draw[thick] (66,0.5)--(66,1.1);
    \draw[thick] (41,0.8)--(56,0.8);
    \draw[thick] (76,0.8)--(94,0.8);
    \draw[thick] (41,0.65)--(41,0.95);
    \draw[thick] (94,0.65)--(94,0.95);
    \node at (22,0.8) {$\times$};
  \end{scope}
\end{tikzpicture}
\end{center}

Box from $56$ to $76$ with the median line at $66$ \hfill \textbf{A1}

Whiskers to $41$ and $94$ \hfill \textbf{A1}

Outlier $22$ marked with a cross \hfill \textbf{A1}

\emph{Award each mark independently; values within $\pm 1$ are accepted. A
lower whisker drawn to $22$ or to the fence $26$ loses the second A1.
FT: a candidate who concluded in (b) that $22$ is not an outlier and draws the
whisker to $22$ with no cross earns the second A1 but not the third.}

\textbf{(d)} Centre: the Class B median ($66$) is higher than the Class A
median ($60$), so Class B generally scored higher. \hfill \textbf{A1}

Spread: the Class B IQR ($20$) is larger than the Class A IQR ($16$), so the
Class B scores are more spread out. \hfill \textbf{A1}

\emph{Each A1 needs a direction stated and the values (or the statistic)
quoted. Accept the ranges, $72$ against $44$, for the spread mark; feedback:
the IQR is the better measure here because the range is driven by the Class B
outlier. A comparison of means earns no mark --- the means are not known.
FT: award the spread mark from the candidate's own Class A IQR in (a).}

\end{enumerate}
\end{document}
